\section{4. Method}

\subsection{4.1 Overview}

ClusterFormer replaces the standard attention operation with four 
components:

1. \textbf{Fast differentiable K-means} over queries ( 4.2).
2. \textbf{Centroid-to-key cross-attention} ( 4.3).
3. \textbf{Soft assignment with learned temperature} ( 4.4).
4. \textbf{Local sliding window residual} ( 4.5).

\textbf{Figure 4.1} shows the pipeline.

\begin{lstlisting}
    +---------------------------------------------+
    |  Input X $\to$ Q, K, V projections (frozen+LoRA) |
    +---------------------+-----------------------+
                          |
              +-----------+-----------+
                                       
        +-----------+           +-----------+
        | K-means   |           | Local     |
        | on Q      |           | Window    |
        | $\to$ C       |           | Attention |
        +-----+-----+           +-----+-----+
              |                       |
                                      |
        +-----------+                 |
        | Centroid- |                 |
        | to-Key    |                 |
        | $\to$ V_cent  |                 |
        +-----+-----+                 |
              |                       |
                                      |
        +-----------+                 |
        | Soft      |                 |
        | Assign    |                 |
        | $\to$ Global  |                 |
        +-----+-----+                 |
              |                       |
              +---------+-------------+
                         
                +--------------+
                | Blend with $\alpha$ |
                +------+-------+
                       |
                        
                  +--------+
                  | out    |
                  +--------+
\end{lstlisting}

\subsection{4.2 Phase 1: Fast Differentiable K-Means}

\subsubsection{4.2.1 Initialization}

Given queries $Q \in \mathbb{R}^{B \times H \times L \times d_h}$, we need 
$K$ initial centroids. Padding tokens must be excluded.

\textbf{Algorithm.} Given a valid-token mask $m \in \{0, 1\}^{B \times L}$:

1. Compute the rank of each valid token: $r_i = \sum_{j \leq i} m_j - 1$.
2. Compute target ranks: $t_k = \lfloor k \cdot N_\text{valid} / K \rfloor$ 
   for $k = 0, \dots, K-1$.
3. For each centroid $k$, pick the valid token whose rank is closest to $t_k$.

This spreads initial centroids uniformly across the valid sequence.

\subsubsection{4.2.2 Lloyd Iterations}

We run \textbf{3 iterations} of standard Lloyd's algorithm:

\textbf{Assignment:}
\[
a_{b,h,i} = \arg\min_{k \in [K]} \|Q_{b,h,i} - C_{b,h,k}\|_2^2
\tag{4.1}
\]

Implemented via \texttt{torch.cdist} on FP32 for numerical stability.

\textbf{Update:}
\[
C_{b,h,k} = \frac{\sum_{i: a_{b,h,i} = k} Q_{b,h,i}}{\sum_{i: a_{b,h,i} = k} 1}
\tag{4.2}
\]

Counts are clamped to $\ge$1 to avoid division by zero. Empty clusters keep 
their previous centroid.

\subsubsection{4.2.3 Why 3 Iterations?}

Empirically, 3 iterations are enough for attention clustering:

\begin{itemize}
\item More iterations increase cost $\mathcal{O}(K L d)$ per iteration.
\item The soft-assignment step (Phase 4) corrects for imperfect centroids.
\item We ablate this in Section 7.6.
\end{itemize}

\subsection{4.3 Phase 2: Centroid-to-Key Cross-Attention}

For each centroid $C_k$, compute attention to all keys:

\[
\text{scores}_{b,h,k,j} = \frac{C_{b,h,k} \cdot K_{b,h,j}}{\sqrt{d_h}}
\tag{4.3}
\]

\[
\alpha_{b,h,k,j} = \text{softmax}_j(\text{scores}_{b,h,k,j})
\tag{4.4}
\]

\[
V_{\text{cent}, b,h,k} = \sum_j \alpha_{b,h,k,j} V_{b,h,j}
\tag{4.5}
\]

\textbf{Complexity.} $\mathcal{O}(K \cdot L \cdot d_h)$ -- linear in $L$!

\textbf{Padding.} Invalid keys are masked with $-\infty$ before softmax.

\textbf{Causality.} For decoders, we apply a causal mask by construction: 
scores are restricted to $j \leq$ some threshold. See Section 5.4.

\subsection{4.4 Phase 3: Soft Assignment with Temperature}

Each query is assigned to centroids via soft attention:

\[
\text{assign}_{b,h,i,k} = \text{softmax}_k\!\left(\frac{Q_{b,h,i} \cdot C_{b,h,k}}{\sqrt{d_h} \cdot \tau}\right)
\tag{4.6}
\]

with learnable temperature $\tau$. The global context is:

\[
\text{Global}_{b,h,i} = \sum_k \text{assign}_{b,h,i,k} \cdot V_{\text{cent}, b,h,k}
\tag{4.7}
\]

\textbf{Why temperature?} Small $\tau$ sharpens the assignment toward the 
nearest centroid (harder clustering); large $\tau$ softens it. We initialize 
$\tau = 0.1$ and let recovery training learn the optimal value.

\textbf{Clamping.} $\tau$ is clamped to $[0.1, 10.0]$ for stability.

\textbf{Empty clusters.} If a cluster has zero assigned tokens, its assignment 
logit is masked to $-\infty$, ensuring no query attends to it.

\subsection{4.5 Phase 4: Local Sliding Window Residual}

Clustering captures global structure but can miss fine-grained local 
dependencies. We add a sliding-window attention branch:

\[
\text{scores}^\text{local}_{b,h,i,j} = \frac{Q_{b,h,i} \cdot K_{b,h,j}}{\sqrt{d_h}} \cdot \mathbf{1}[|i - j| \leq W]
\tag{4.8}
\]

\[
\text{Local}_{b,h,i} = \sum_{j: |i-j| \leq W} \text{softmax}_j(\text{scores}^\text{local}) \cdot V_{b,h,j}
\tag{4.9}
\]

\textbf{Complexity.} $\mathcal{O}(L \cdot W \cdot d_h)$.

\subsubsection{4.5.1 Blending}

The final output is:

\[
\text{Output} = (1 - \sigma(\alpha)) \cdot \text{Global} + \sigma(\alpha) \cdot \text{Local}
\tag{4.10}
\]

where $\alpha \in \mathbb{R}$ is a \textbf{learnable scalar} initialized to 0 
(so $\sigma(0) = 0.5$, giving 50/50 blend at start).

\textbf{Why sigmoid?} Enforces $\alpha_\text{blend} \in (0, 1)$, avoiding 
negative or >1 weights.

\subsection{4.6 The Full Algorithm}

\textbf{Algorithm 1:} ClusterFormer forward pass
\begin{lstlisting}
Input:  X $\in$ R^{B$\times$L$\times$d}, attention_mask, position_embeddings
Output: Y, attn_weights

 1: Q $\leftarrow$ LoRALinear_Q(X)                         Frozen base + LoRA
 2: K $\leftarrow$ LoRALinear_K(X)
 3: V $\leftarrow$ LoRALinear_V(X)
 4:
 5: (Q, K) $\leftarrow$ apply_rope(Q, K, cos, sin)
 6: if GQA: K, V $\leftarrow$ repeat_kv(K, V)
 7:
 8: if decoder:                                   Exact causal path
 9:     S $\leftarrow$ QK^T / $\sqrt{}$d_h
10:     S $\leftarrow$ causal_mask(S)
11:     A $\leftarrow$ softmax(S)
12:     Y_ctx $\leftarrow$ A @ V
13: else:                                          Cluster path
14:     C, ids, counts $\leftarrow$ kmeans(Q, K, iters=3)     Cluster queries
15:     S_cent $\leftarrow$ C K^T / $\sqrt{}$d_h
16:     A_cent $\leftarrow$ softmax(S_cent)
17:     V_cent $\leftarrow$ A_cent @ V                         Centroid contexts
18:     S_assign $\leftarrow$ Q C^T / ($\sqrt{}$d_h   $\tau$)
19:     A_assign $\leftarrow$ softmax(S_assign)
20:     Y_global $\leftarrow$ A_assign @ V_cent
21:
22:     if W > 0:
23:         S_local $\leftarrow$ QK^T / $\sqrt{}$d_h
24:         S_local $\leftarrow$ window_mask(S_local, W)
25:         A_local $\leftarrow$ softmax(S_local)
26:         Y_local $\leftarrow$ A_local @ V
27:         $\alpha$ $\leftarrow$ $\sigma$(learned_scalar)
28:         Y_ctx $\leftarrow$ (1-$\alpha$) Y_global + $\alpha$ Y_local
29:     else:
30:         Y_ctx $\leftarrow$ Y_global
31:
32: Y $\leftarrow$ Y_ctx.reshape(B, L, H d_h)
33: if nested: Y_out $\leftarrow$ Y
34: else: Y_out $\leftarrow$ LoRALinear_O(Y)
35:
36: return Y_out
\end{lstlisting}

\subsection{4.7 Complexity Analysis}

\subsubsection{4.7.1 Time Complexity}

\begin{table}[H]
\centering
\small
\begin{tabular}{ll}
\toprule
\textbf{Component} & \textbf{Complexity} \\
\midrule
Q/K/V projections & $\mathcal{O}(L d^2)$ \\
K-means (3 iters) & $\mathcal{O}(K L d_h)$ \\
Centroid-to-key & $\mathcal{O}(K L d_h)$ \\
Soft assignment & $\mathcal{O}(K L d_h)$ \\
Local window & $\mathcal{O}(L W d_h)$ \\
Output projection & $\mathcal{O}(L d^2)$ \\
**Total** & $\mathcal{O}(L d^2 + L K d_h + L W d_h)$ \\
\bottomrule
\end{tabular}
\end{table}

For $K, W \ll L$, this is effectively \textbf{linear in $L$}.

\subsubsection{4.7.2 Comparison with Baseline}

\begin{table}[H]
\centering
\small
\begin{tabular}{ll}
\toprule
\textbf{Config} & \textbf{Attention cost} \\
\midrule
Baseline & $\mathcal{O}(L^2 d_h)$ \\
ClusterFormer ($K=16, W=32$) & $\mathcal{O}(L \cdot 48 \cdot d_h)$ \\
**Reduction at $L=128$** & **$\approx 2.7\times$** \\
**Reduction at $L=512$** & **$\approx 10.7\times$** \\
\bottomrule
\end{tabular}
\end{table}

\subsubsection{4.7.3 Parameter Count}

\begin{table}[H]
\centering
\small
\begin{tabular}{ll}
\toprule
\textbf{Component} & \textbf{Parameters} \\
\midrule
LoRA ($r_L=16$) & $4 \times r_L \times d$ per layer \\
Learnable temperature & 1 \\
Learnable blend & 1 \\
**Per layer** & $4 r_L d + 2$ \\
**Total ($N_l$ layers)** & $4 r_L d N_l + 2 N_l$ \\
\bottomrule
\end{tabular}
\end{table}

For DistilBERT ($d=768$, $r_L=16$, $N_l=6$): \textbf{+294,916 params}.

\textbf{ClusterFormer is not a compression method.} Parameter count increases 
by \textbf{0.88\%}.

\subsection{4.8 Theoretical Notes}

\subsubsection{4.8.1 K-Means Convergence}

Lloyd's algorithm monotonically decreases the objective (3.3) and converges 
in finite iterations. We use 3 iterations as a practical approximation.

\subsubsection{4.8.2 Complexity Proof}

\textbf{Claim.} ClusterFormer's attention cost is $\mathcal{O}(K L d_h)$.

\textbf{Proof.} Direct from the algorithm:

\begin{itemize}
\item K-means: each iteration is $\mathcal{O}(L K d_h)$ via cdist. Three
\end{itemize}
  iterations = $\mathcal{O}(3 L K d_h)$.
\begin{itemize}
\item Centroid-to-key: matrix product of $(K \times d_h)$ and $(d_h \times L)$
\end{itemize}
  is $\mathcal{O}(K L d_h)$.
\begin{itemize}
\item Soft assignment: product of $(L \times d_h)$ and $(d_h \times K)$ is
\end{itemize}
  $\mathcal{O}(L K d_h)$.

Sum is $\mathcal{O}(L K d_h)$. $\square$

\subsubsection{4.8.3 Approximation vs Exactness}

Unlike Performer/Linformer, ClusterFormer is \textbf{exact} for the attention 
distribution over centroids. It approximates the attention over tokens by 
grouping tokens. The approximation error is bounded by the K-means 
distortion:

\[
\epsilon \leq \sum_i \|Q_i - C_{a_i}\|_2^2
\tag{4.11}
\]

See Appendix A.5 for details.

\subsection{4.9 Hyperparameter Guidance}

\begin{table}[H]
\centering
\small
\begin{tabular}{lll}
\toprule
\textbf{Hyperparameter} & \textbf{Default} & \textbf{Effect} \\
\midrule
$K$ (num\_clusters) & 16 & Higher = better fidelity, more cost \\
Iterations & 3 & Higher = better centroids \\
$W$ (window\_size) & 32 & Higher = more local context \\
$\tau$ (temp) & learnable & Smaller = harder clustering \\
$\alpha$ (blend) & learnable & Higher = more local \\
\bottomrule
\end{tabular}
\end{table}

\textbf{Choosing K:}

\begin{itemize}
\item $K = 8$: aggressive clustering, faster
\item $K = 16$: balanced (default)
\item $K = 32$: more fidelity, slower
\item $K = L$: no compression (full attention)
\end{itemize}